% Source: https://artofproblemsolving.com/community/c3h2368783_geometry_problem_4
% Stable identity: aops-topic-2368783-q04
\textbf{Setting.} Let $O$ be the centre. Since the square is convex, the random points are the vertices, in order around it, of a convex polygon, namely their convex hull. (If several points lie on the same side of the square, some angles of this polygon are straight angles; this does not matter, since only the region it encloses is used.) Ignore events of probability $0$ (two points equal, or two points on a line through $O$).

\textbf{Step 1.} $O$ lies outside the polygon exactly when all the points lie in an open half-plane bounded by a line through $O$. Indeed, if all points lie in such a half-plane, so does their convex hull; and if $O$ is outside the convex polygon, a line separating $O$ from the polygon, moved parallel to itself until it passes through $O$, leaves the polygon strictly on one side.

\textbf{Step 2.} For each point $X_i$, let $H_i$ be the open half-plane obtained by turning the ray $OX_i$ counterclockwise through angles strictly between $0^\circ$ and $180^\circ$, and let $E_i$ be the event that all the other points lie in $H_i$. If all points lie in an open half-plane through $O$, that half-plane is swept out by a ray from $O$ turning counterclockwise through $180^\circ$ from one of its bounding rays; let $X_i$ be the first point the ray meets. All the other points are met after turning less than $180^\circ$ further, so $E_i$ holds. Conversely, if $E_i$ holds, turning the line $OX_i$ slightly clockwise gives an open half-plane containing all the points. And two of the events cannot both hold: $X_j\in H_i$ and $X_i\in H_j$ would mean that going counterclockwise from $X_i$ to $X_j$ and from $X_j$ back to $X_i$ both take less than $180^\circ$, although together they make $360^\circ$.

\textbf{Step 3.} The reflection in $O$ maps the boundary of the square onto itself and preserves the uniform distribution on it, and it swaps $H_i$ with the opposite open half-plane. So, whatever $X_i$ is, each other point lies in $H_i$ with probability $\dfrac12$, independently of the others, and $P(E_i)=\left(\dfrac12\right)^{m-1}$, where $m$ is the number of points.

\textbf{Conclusion.} Hence
\[P(O\text{ outside the polygon})=P(E_1)+\cdots+P(E_m)=\dfrac{m}{2^{m-1}}.\]

\begin{asy}
usepackage("amsmath");
size(7cm);
defaultpen(fontsize(10pt));
void pic(pair c, real[] t, bool inside){
  pair P(real s){ if(s<2) return c+(1,-1+s); if(s<4) return c+(1-(s-2),1); if(s<6) return c+(-1,1-(s-4)); return c+(-1+(s-6),-1);}
  draw(shift(c)*box((-1,-1),(1,1)),linewidth(0.8));
  pair[] p; for(int i=0;i<t.length;++i) p.push(P(t[i]));
  path g=p[0]; for(int i=1;i<t.length;++i) g=g--p[i]; g=g--cycle;
  fill(g,inside?rgb(0.8,0.9,1):rgb(1,0.88,0.85)); draw(g,inside?blue:red);
  for(int i=0;i<t.length;++i) dot(p[i]);
  dot(c); label("$O$",c,SE,fontsize(8pt));
}
pic((0,0),new real[]{0.5,2.4,4.1,5.6,7.0},true);
pic((3.2,0),new real[]{6.4,7.1,7.8,0.9,1.6},false);
draw((3.2,0)+1.3*dir(45)--(3.2,0)+1.3*dir(225),dashed);
label("contains $O$",(0,-1.25),S,fontsize(8pt));
label("in a half-plane",(3.2,-1.25),S,fontsize(8pt));
\end{asy}

With $m=5$ points, $P(O\text{ outside})=\dfrac{5}{16}$, so the probability that the centre lies inside the pentagon is $\boxed{\dfrac{11}{16}}$.
